\chapter{符号与缩写}\label{app:symbols}

\subsection*{字符}
\begin{center}
\small
\renewcommand{\arraystretch}{1.15}
\begin{tabular}{>{\centering\arraybackslash}p{0.16\textwidth}p{0.48\textwidth}>{\centering\arraybackslash}p{0.18\textwidth}}
\hline
\textbf{Symbol} & \textbf{Description} & \textbf{Unit} \\
\hline
$\mathbf{p}$ & 位置向量 & $\Unit{m}$ \\
$\mathbf{v}$ & 线速度向量 & $\Unit{m \cdot s^{-1}}$ \\
$\boldsymbol{\omega}$ & 角速度向量 & $\Unit{rad \cdot s^{-1}}$ \\
$\mathbf{q}$ & 单位四元数 & $\Unit{1}$ \\
$\Delta t$ & 仿真步长 & $\Unit{s}$ \\
$a_t$ & 策略在时刻$t$输出的任务空间动作 & $\Unit{1}$ \\
$o_t$ & 策略在时刻$t$接收的观测 & $\Unit{1}$ \\
$O_t$ & 历史观测窗口 & $\Unit{1}$ \\
$C_t$ & 未来目标条件窗口 & $\Unit{1}$ \\
$A_t$ & 动作块序列 & $\Unit{1}$ \\
$H_o$ & 历史观测窗口长度 & $\Unit{1}$ \\
$H_a$ & 动作块预测长度 & $\Unit{1}$ \\
$H_c$ & 未来条件窗口长度 & $\Unit{1}$ \\
$\phi_t$ & 离散任务阶段索引 & $\Unit{1}$ \\
$\eta_t$ & 观测噪声项 & $\Unit{1}$ \\
\hline
\end{tabular}
\end{center}

\subsection*{算子}
\begin{center}
\small
\renewcommand{\arraystretch}{1.15}
\begin{tabular}{>{\centering\arraybackslash}p{0.22\textwidth}p{0.45\textwidth}}
\hline
\textbf{Symbol} & \textbf{Description} \\
\hline
$\Delta$ & 差分或相对量 \\
$\|\cdot\|$ & 欧氏范数 \\
$\mathbb{E}[\cdot]$ & 期望 \\
$\otimes$ & 四元数乘法 \\
$R(q)$ & 四元数$q$对应的旋转矩阵 \\
$d_R(\cdot,\cdot)$ & 四元数姿态测地距离 \\
\hline
\end{tabular}
\end{center}

\subsection*{缩写}
\begin{center}
\small
\renewcommand{\arraystretch}{1.12}
\begin{tabular}{>{\centering\arraybackslash}p{0.22\textwidth}p{0.52\textwidth}}
\hline
\textbf{Abbreviation} & \textbf{Description} \\
\hline
ACT & Action Chunking Transformer, 动作分块变换器 \\
BC & Behavior Cloning, 行为克隆 \\
BC-RNN & Recurrent Behavior Cloning, 循环行为克隆 \\
CVAE & Conditional Variational Autoencoder, 条件变分自编码器 \\
BCE & Binary Cross Entropy, 二元交叉熵 \\
CE & Cross Entropy, 交叉熵 \\
DDIM & Denoising Diffusion Implicit Models, 去噪扩散隐式模型 \\
DP & Diffusion Policy, 扩散策略 \\
EE & End Effector, 末端执行器 \\
EMA & Exponential Moving Average, 指数滑动平均 \\
FFN & Feed-Forward Network, 前馈网络 \\
FSM & Finite-State Machine, 有限状态机 \\
GT & Ground Truth, 仿真真值观测 \\
IK & Inverse Kinematics, 逆运动学 \\
IK+EMA & Inverse Kinematics with Exponential Moving Average, 逆运动学与指数滑动平均跟踪执行器 \\
KL & Kullback-Leibler Divergence, KL散度 \\
MPC & Model Predictive Control, 模型预测控制 \\
OTF-noise & On-the-Fly Observation Noise, 在线观测噪声增强训练 \\
SE(3) & Special Euclidean Group, 三维刚体位姿群 \\
SR & Success Rate, 闭环成功率 \\
\hline
\end{tabular}
\end{center}


\chapter{实验复现与完整结果}\label{app:reproducibility}

\section{实验复现信息}

闭环成功率定义为成功回合数占总评估回合数的比例。除特别说明外，每个策略-环境组合包含750个闭环评估回合，所有成功率均附带95\% Wilson置信区间。

\begin{table}[!htbp]
    \caption{实验环境配置摘要。}
    \label{tab:env-summary}
    \centering
    \small
    \renewcommand{\arraystretch}{1.15}
    \begin{tabularx}{\textwidth}{>{\raggedright\arraybackslash}p{0.20\textwidth}>{\raggedright\arraybackslash}p{0.16\textwidth}>{\raggedright\arraybackslash}p{0.13\textwidth}Y}
        \hline
        场景 & 评估配置 & 观测流 & 关键运动设置 \\
        \hline
        翻滚平移，无噪观测 & 无噪观测配置 & 真实状态观测 & 随机初始姿态，空间轴$15\,\mathrm{deg/s}$旋转，线速度范围$[-0.075,-0.075,-0.02]$到$[0.075,0.075,0.02]$ \\
        翻滚平移，噪声观测 & 噪声观测配置 & 含噪观测流 & 与翻滚平移无噪观测相同运动，加入漂移、高斯、离群点、随机丢弃和2帧时延 \\
        自旋平移，无噪观测 & 无噪观测配置 & 真实状态观测 & 固定姿态，绕$z$轴$30$--$60\,\mathrm{deg/s}$自旋，线速度范围$[-0.10,-0.10,-0.02]$到$[0.10,0.10,0.02]$ \\
        自旋平移，噪声观测 & 噪声观测配置 & 含噪观测流 & 与自旋平移无噪观测相同运动，加入漂移、高斯噪声、离群点、随机丢弃和2帧时延 \\
        \hline
    \end{tabularx}
\end{table}

主数据集使用完整目标中心表示（历史长度16），从5000条成功专家示教中构造。随机姿态翻滚平移场景包含约19.4万个训练窗口和约1.0万个验证窗口，定轴自旋平移场景包含约19.9万个训练窗口和约1.0万个验证窗口。

数据处理规则：(1) 动作标签取自专家参考轨迹；(2) 仅保留抓取相关阶段的窗口；(3) 采样步幅为4；(4) 验证比例为0.05；(5) 数据划分使用固定随机种子以保证可复现性。统一采用动作块长度32、未来条件长度50、仿真步长0.01 s。

数据量消融中的10\%、25\%、50\%和100\%分别对应500、1250、2500和5000条成功示教；除示教条数外，数据处理规则保持一致。

\begin{table}[!htbp]
    \caption{训练数据构造明细。窗口数统计对应完整目标中心表示、历史长度16和100\%示教数据。}
    \label{tab:dataset-construction-details}
    \centering
    \small
    \renewcommand{\arraystretch}{1.15}
    \begin{tabularx}{\textwidth}{>{\raggedright\arraybackslash}p{0.22\textwidth}>{\centering\arraybackslash}p{0.13\textwidth}>{\centering\arraybackslash}p{0.20\textwidth}Y}
        \hline
        场景 & 成功示教数 & 训练/验证窗口数 & 数据处理说明 \\
        \hline
        随机姿态翻滚平移 & 5000 & 193575 / 10203 & 初始姿态随机，空间轴旋转；仅保留抓取相关阶段窗口，动作标签来自专家参考轨迹。 \\
        定轴自旋平移 & 5000 & 198936 / 9921 & 初始姿态固定，绕竖直轴自旋；窗口生成、验证划分和动作标签规则与翻滚场景一致。 \\
        \hline
    \end{tabularx}
\end{table}

\section{训练与运行时超参数}

所有策略均训练1500轮，共享以下训练协议：AdamW优化器、余弦退火学习率调度、0.999的EMA衰减系数和1.0的梯度范数裁剪阈值。

统一动作接口采用长度为32的动作块，闭环执行时每次提交前16步动作，闭合预测阈值为0.5。在翻滚平移与自旋平移两个场景中，BC、BC-RNN和ACT使用相同的优化器协议，而DP 12M在两个场景中采用了不同的学习率和预热步数。表~\ref{tab:training-runtime-protocol}列出各策略存在差异的训练参数和模型选择指标。

\begin{table}[!htbp]
    \caption{主策略差异训练参数。权重衰减、批次大小和预热步数随策略而变，未列出的运行时协议在各策略间保持一致。}
    \label{tab:training-runtime-protocol}
    \centering
    \small
    \renewcommand{\arraystretch}{1.15}
    \begin{tabularx}{\textwidth}{>{\raggedright\arraybackslash}p{0.21\textwidth}>{\centering\arraybackslash}p{0.13\textwidth}>{\centering\arraybackslash}p{0.13\textwidth}>{\centering\arraybackslash}p{0.13\textwidth}>{\centering\arraybackslash}p{0.11\textwidth}Y}
        \hline
        策略/场景 & 学习率 & 批次大小 & 权重衰减 & 预热步数 & 模型选择指标 \\
        \hline
        BC & $2\times10^{-4}$ & 16384 & $1\times10^{-6}$ & 300 & 验证动作损失最低 \\
        BC-RNN & $2\times10^{-4}$ & 8192 & $1\times10^{-6}$ & 300 & 验证动作损失最低 \\
        ACT & $2\times10^{-4}$ & 1024 & $1\times10^{-4}$ & 1000 & 验证动作损失最低 \\
        DP 12M（翻滚） & $1\times10^{-4}$ & 12288 & $1\times10^{-6}$ & 500 & 验证噪声预测损失最低 \\
        DP 12M（自旋） & $2\times10^{-4}$ & 12288 & $1\times10^{-6}$ & 300 & 验证噪声预测损失最低 \\
        \hline
    \end{tabularx}
\end{table}

训练目标由主任务损失和辅助监督组成。BC与BC-RNN使用动作$L_1$损失作为主指标；ACT在动作损失之外加入CVAE的KL正则项；DP使用扩散噪声预测损失作为主指标。所有策略均包含闭合预测头和阶段预测头，闭合预测头使用二分类BCE损失，阶段预测头使用5类CE损失。表~\ref{tab:loss-selection-details}给出损失权重和早停规则。

\begin{table}[!htbp]
    \caption{训练损失和模型选择规则。}
    \label{tab:loss-selection-details}
    \centering
    \small
    \renewcommand{\arraystretch}{1.15}
    \begin{tabularx}{\textwidth}{>{\raggedright\arraybackslash}p{0.15\textwidth}YYY}
        \hline
        模型 & 主损失与权重 & 辅助监督与正则项 & 模型选择与早停规则 \\
        \hline
        BC & 动作$L_1$，权重1.0 & 闭合BCE权重5.0；阶段CE权重1.0 & 选择验证动作损失最低时的模型参数；验证损失连续10轮未下降则触发早停。 \\
        BC-RNN & 动作$L_1$，权重1.0 & 闭合BCE权重5.0；阶段CE权重1.0 & 选择验证动作损失最低时的模型参数；验证损失连续10轮未下降则触发早停。 \\
        ACT & 动作$L_1$，权重1.0 & KL权重最高10.0，前10000步线性预热；闭合BCE权重5.0；阶段CE权重1.0 & 选择验证动作损失最低时的模型参数；验证损失连续5轮未下降则触发早停。 \\
        DP & 噪声预测损失，权重1.0 & 闭合BCE权重5.0；阶段CE权重1.0 & 选择验证噪声预测损失最低时的模型参数；不使用早停，完整训练1500轮。 \\
        \hline
    \end{tabularx}
\end{table}

\begin{table}[!htbp]
    \caption{定轴自旋平移场景新增实验设置明细。表中列出每组实验改变的因素；其余训练协议遵循表~\ref{tab:training-runtime-protocol}和表~\ref{tab:loss-selection-details}。}
    \label{tab:spin-training-details}
    \centering
    \small
    \renewcommand{\arraystretch}{1.15}
    \begin{tabularx}{\textwidth}{>{\raggedright\arraybackslash}p{0.16\textwidth}>{\raggedright\arraybackslash}p{0.21\textwidth}YY}
        \hline
        实验族 & 涉及模型 & 改变因素 & 固定设置 \\
        \hline
        主策略对比 & BC、BC-RNN、ACT、DP 12M & 策略架构 & 完整目标中心表示，$H_o=16$，无噪声训练数据；训练集198936个窗口，验证集9921个窗口。 \\
        表示消融 & BC、ACT & 坐标系和历史表示维度 & $H_o=16$；动作块长度、执行窗口、优化器协议和模型选择规则保持不变。 \\
        噪声增强 & BC-RNN、ACT、DP 12M & 训练时观测扰动 & 完整目标中心表示，$H_o=16$；评估时使用含噪观测流。 \\
        数据量消融 & BC-RNN、ACT & 成功示教子集规模 & 完整目标中心表示，$H_o=16$；验证比例、采样步幅、动作块长度和未来条件长度保持一致。 \\
        历史窗口长度消融 & BC-RNN、ACT & 历史窗口长度$H_o$ & 紧凑目标中心表示；动作块长度为32，未来条件长度为50，执行窗口为16。 \\
        容量与DP参考 & BC-RNN 6M、DP 1.5M/6M/12M/18M & 模型容量、坐标表示和噪声设置 & 完整目标中心表示，$H_o=16$；DP容量变体使用相同扩散步数和DDIM推理步数。 \\
        \hline
    \end{tabularx}
\end{table}

表示消融中的四种历史输入分别为目标中心紧凑表示、目标中心完整表示、世界坐标紧凑表示和世界坐标完整表示，对应历史输入维度8、14、15和21。未来条件维度均为7，动作维度均为8。

噪声增强实验在训练时在线扰动观测序列。注入噪声的参数如下：位置噪声标准差为0.03 m，姿态噪声标准差为$10^\circ$，线速度噪声标准差为0.02 m/s，角速度噪声标准差为$5^\circ$/s；此外设置随机丢弃率为0.05、漂移标准差为0.01、离群点率为0.02、离群幅值为0.15、观测时延为2帧。

容量实验中，BC-RNN 6M将GRU隐藏维度从218增至473（对应参数量从约1.5M增至约6M）。DP容量变体在自旋场景中的学习率均为$2\times10^{-4}$，预热步数均为300。

\begin{table}[!htbp]
    \caption{定轴自旋平移场景主策略、噪声增强和容量实验的最低验证损失。BC、BC-RNN和ACT列出验证动作损失；DP列出验证噪声预测损失。}
    \label{tab:spin-best-validation-loss}
    \centering
    \small
    \renewcommand{\arraystretch}{1.12}
    \begin{tabularx}{\textwidth}{>{\raggedright\arraybackslash}p{0.18\textwidth}Y>{\raggedright\arraybackslash}p{0.20\textwidth}>{\centering\arraybackslash}p{0.18\textwidth}}
        \hline
        实验族 & 设置 & 模型 & 最低验证损失 \\
        \hline
        主策略 & 完整目标中心，$H_o=16$ & BC & 0.172 \\
        主策略 & 完整目标中心，$H_o=16$ & BC-RNN & 0.125 \\
        主策略 & 完整目标中心，$H_o=16$ & ACT & 0.103 \\
        主策略 & 完整目标中心，$H_o=16$ & DP 12M & 0.022 \\
        噪声增强 & 完整目标中心，$H_o=16$ & BC-RNN & 0.149 \\
        噪声增强 & 完整目标中心，$H_o=16$ & ACT & 0.147 \\
        噪声增强 & 完整目标中心，$H_o=16$ & DP 12M & 0.031 \\
        容量参考 & 完整目标中心，$H_o=16$ & BC-RNN 6M & 0.145 \\
        容量参考 & 完整目标中心，$H_o=16$ & DP 1.5M & 0.071 \\
        容量参考 & 完整目标中心，$H_o=16$ & DP 6M & 0.027 \\
        容量参考 & 完整目标中心，$H_o=16$ & DP 12M & 0.022 \\
        容量参考 & 完整目标中心，$H_o=16$ & DP 18M & 0.021 \\
        \hline
    \end{tabularx}
\end{table}

\begin{table}[!htbp]
    \caption{定轴自旋平移场景数据规模消融的最低验证动作损失。}
    \label{tab:spin-data-scale-validation-loss}
    \centering
    \small
    \renewcommand{\arraystretch}{1.12}
    \begin{tabularx}{\textwidth}{>{\raggedright\arraybackslash}p{0.18\textwidth}Y>{\raggedright\arraybackslash}p{0.18\textwidth}>{\centering\arraybackslash}p{0.18\textwidth}}
        \hline
        模型 & 固定设置 & 数据比例 & 最低验证损失 \\
        \hline
        BC-RNN & 完整目标中心表示，$H_o=16$ & 10\% & 0.260 \\
        BC-RNN & 完整目标中心表示，$H_o=16$ & 25\% & 0.164 \\
        BC-RNN & 完整目标中心表示，$H_o=16$ & 50\% & 0.138 \\
        BC-RNN & 完整目标中心表示，$H_o=16$ & 100\% & 0.125 \\
        ACT & 完整目标中心表示，$H_o=16$ & 10\% & 0.136 \\
        ACT & 完整目标中心表示，$H_o=16$ & 25\% & 0.113 \\
        ACT & 完整目标中心表示，$H_o=16$ & 50\% & 0.105 \\
        ACT & 完整目标中心表示，$H_o=16$ & 100\% & 0.103 \\
        \hline
    \end{tabularx}
\end{table}

\begin{table}[!htbp]
    \caption{定轴自旋平移场景历史窗口长度消融的最低验证动作损失。}
    \label{tab:spin-history-validation-loss}
    \centering
    \small
    \renewcommand{\arraystretch}{1.12}
    \begin{tabularx}{\textwidth}{>{\raggedright\arraybackslash}p{0.18\textwidth}Y>{\raggedright\arraybackslash}p{0.18\textwidth}>{\centering\arraybackslash}p{0.18\textwidth}}
        \hline
        模型 & 固定设置 & 历史窗口长度 & 最低验证损失 \\
        \hline
        BC-RNN & 紧凑目标中心表示 & $H_o=1$ & 0.132 \\
        BC-RNN & 紧凑目标中心表示 & $H_o=8$ & 0.124 \\
        BC-RNN & 紧凑目标中心表示 & $H_o=32$ & 0.116 \\
        BC-RNN & 紧凑目标中心表示 & $H_o=64$ & 0.122 \\
        ACT & 紧凑目标中心表示 & $H_o=1$ & 0.128 \\
        ACT & 紧凑目标中心表示 & $H_o=8$ & 0.108 \\
        ACT & 紧凑目标中心表示 & $H_o=32$ & 0.103 \\
        ACT & 紧凑目标中心表示 & $H_o=64$ & 0.101 \\
        \hline
    \end{tabularx}
\end{table}

\begin{table}[!htbp]
    \caption{BC、BC-RNN和ACT结构补充参数。所有模型均使用闭合预测头和5类阶段预测头，动作块长度为32。}
    \label{tab:model-structure-details}
    \centering
    \small
    \renewcommand{\arraystretch}{1.15}
    \begin{tabularx}{\textwidth}{>{\raggedright\arraybackslash}p{0.14\textwidth}YY>{\centering\arraybackslash}p{0.13\textwidth}>{\centering\arraybackslash}p{0.13\textwidth}>{\centering\arraybackslash}p{0.11\textwidth}}
        \hline
        模型 & 编码结构 & 主体结构 & 历史维度 & 未来维度 & 动作维度 \\
        \hline
        BC & 历史编码通道为64、128、192，卷积核5；未来条件编码通道为64、128；条件维度192。 & 3层MLP，每层480维。 & 14 & 7 & 8 \\
        BC-RNN & 未来条件编码通道$(64,128)$；条件维度192。 & 2层GRU，隐藏维度218，随机失活率0.1；6M扩容版本隐藏维度473。 & 14 & 7 & 8 \\
        ACT & CVAE编码器2层，隐变量维度32。 & Transformer维度132，4个注意力头，2层编码器，2层解码器，FFN维度544，随机失活率0.1。 & 14 & 7 & 8 \\
        \hline
    \end{tabularx}
\end{table}

\begin{table}[!htbp]
    \caption{DP容量变体结构补充参数。所有DP变体均使用条件1D U-Net，100个训练扩散步和10个推理采样步；历史输入维度为14，未来条件维度为7，动作维度为8。}
    \label{tab:dp-structure-details}
    \centering
    \small
    \renewcommand{\arraystretch}{1.15}
    \begin{tabularx}{\textwidth}{>{\raggedright\arraybackslash}p{0.16\textwidth}YY>{\centering\arraybackslash}p{0.18\textwidth}>{\centering\arraybackslash}p{0.13\textwidth}}
        \hline
        容量变体 & 编码通道 & U-Net设置 & 通道倍数 & 批次大小 \\
        \hline
        DP 1.5M & 历史通道为64、128、128；未来通道为64、128。 & 条件维度80；基通道64。 & 1、2、2 & 32768 \\
        DP 6M & 历史通道为64、128、128；未来通道为64、128。 & 条件维度128；基通道88。 & 1、2、4 & 16384 \\
        DP 12M & 历史通道为32、64、96；未来通道为32、64。 & 条件维度160；基通道128。 & 1、2、4 & 12288 \\
        DP 18M & 历史通道为32、64、96；未来通道为32、64。 & 条件维度160；基通道160。 & 1、2、4 & 12288 \\
        \hline
    \end{tabularx}
\end{table}

\section{完整结果表}

\begin{table}[!htbp]
    \caption{主无噪比较完整表：随机姿态翻滚平移场景。}
    \label{tab:appendix-main-clean}
    \centering
    \small
    \renewcommand{\arraystretch}{1.12}
    \begin{tabularx}{\textwidth}{>{\raggedright\arraybackslash}p{0.15\textwidth}ZZZZZZ}
        \hline
        场景 & BC & BC-RNN & ACT & DP 1.5M & DP 6M & DP 12M \\
        \hline
        翻滚平移，无噪观测 & 0.1\% & 48.7\% & 59.2\% & 1.3\% & 2.8\% & 2.5\% \\
        \hline
    \end{tabularx}
\end{table}


\section{扩散策略容量扩展实验}\label{sec:appendix-dp}

正文§4.1给出了DP 12M在翻滚平移场景中的成功率（2.5\%）。表~\ref{tab:appendix-dp-scaling}进一步列出了DP容量（1.5M至18M）和表示变体在该场景中的完整比较结果，所有结果均基于750回合评估。

该表同时列出定轴自旋平移场景中的DP容量扩展结果，以供对照。DP 12M在该场景中成功率为41.9\%。作为参考，BC在同一场景中成功率为45.7\%，二者处于相近水平。这一结果表明DP在抓取面已知的条件下可以正常工作。相比之下，所有DP变体在翻滚场景中仅达到1.3\%--2.8\%，且置信区间高度重叠，表明容量或表示的变化均无法在变动抓取面的场景中改善DP的闭环控制。由于DP在翻滚场景中所有已评估变体成功率均不超过2.8\%且置信区间高度重叠，未进一步评估18M及世界坐标变体在该场景中的表现，表中对应位置以"—"标出。

\begin{table}[!htbp]
    \caption{DP扩展性和表示结果。验证损失为所选模型的验证集噪声预测损失；方括号内为95\% Wilson置信区间。}
    \label{tab:appendix-dp-scaling}
    \centering
    \small
    \renewcommand{\arraystretch}{1.12}
    \begin{tabularx}{\textwidth}{>{\raggedright\arraybackslash}p{0.30\textwidth}>{\centering\arraybackslash}p{0.14\textwidth}ZZ}
        \hline
        实验 & 验证损失 & 翻滚平移成功率 & 自旋平移成功率（验证参考） \\
        \hline
        DP 1.5M，完整目标中心 & 0.071 & \ResultCI{1.3\%}{[0.7, 2.4]} & \ResultCI{14.9\%}{[12.6, 17.7]} \\
        DP 6M，完整目标中心 & 0.027 & \ResultCI{2.8\%}{[1.8, 4.2]} & \ResultCI{36.9\%}{[33.6, 40.4]} \\
        DP 12M，完整目标中心 & 0.022 & \ResultCI{2.5\%}{[1.6, 3.9]} & \BestResultCI{41.9\%}{[38.4, 45.4]} \\
        DP 18M，完整目标中心 & \textbf{0.021} & -- & \ResultCI{40.4\%}{[36.9, 44.0]} \\
        DP 12M，完整世界坐标 & 0.026 & -- & \ResultCI{2.9\%}{[1.9, 4.4]} \\
        DP 12M，完整世界坐标，噪声增强 & 0.052 & -- & \ResultCI{1.2\%}{[0.6, 2.3]} \\
        \hline
    \end{tabularx}
\end{table}

DP的验证噪声预测损失单调递减（0.071→0.021），但在翻滚平移场景中所有容量变体的成功率均处于1.3\%--2.8\%之间（各置信区间高度重叠），容量增加并未改善闭环成功率。

在定轴自旋平移验证场景中，DP的成功率随容量增加从14.9\%提升至41.9\%，但18M（成功率40.4\%）与12M（成功率41.9\%）的置信区间几乎完全重叠，差异未达统计显著性。可见DP的验证噪声预测损失与闭环成功率之间缺乏可靠对应关系：更低的验证损失并不保证更高的闭环成功率。

使用完整世界坐标表示时，DP 12M仅2.9\%，噪声增强版本进一步降至1.2\%。综上，在本实验条件下DP于12M处容量已近饱和，容量和表示的改变均无法改善DP在翻滚平移场景中的闭环控制表现。

\section{统计口径与图表说明}

正文与附录中的闭环成功率统一定义如下：若回合结束时立方体被稳定运输至目标区域，则该回合记为成功；点估计为成功回合数除以总回合数，区间估计为95\% Wilson置信区间。训练曲线展示每个训练轮次记录的训练损失和验证损失；BC、BC-RNN和ACT绘制动作$L_1$损失，DP绘制噪声预测损失。星号标记的是最低验证损失对应的训练轮次。

\chapter{训练损失曲线}\label{app:loss-curves}

本附录给出正文和附录中涉及的各学习策略训练曲线。BC、BC-RNN和ACT按验证动作损失选择模型参数，DP按验证噪声预测损失选择模型参数；星号标注对应曲线上的最低验证损失。训练损失用于训练监控，正文结论仍以闭环成功率为最终依据。

\begin{figure}[!htbp]
    \centering
    \includegraphics[width=\textwidth]{fig_app_loss_main_policies}
    \caption{主随机姿态翻滚平移实验四类策略的训练损失曲线。前三个子图为动作$L_1$损失，DP子图为噪声预测损失；灰线为训练损失，彩色线为验证损失。}
    \label{fig:app-loss-main-policies}
\end{figure}

\begin{figure}[!htbp]
    \centering
    \includegraphics[width=\textwidth]{fig_app_loss_spin_main_policies}
    \caption{定轴自旋平移验证场景四类主策略的训练损失曲线。前三个子图为动作$L_1$损失，DP子图为噪声预测损失；灰线为训练损失，彩色线为验证损失。}
    \label{fig:app-loss-spin-main-policies}
\end{figure}

\begin{figure}[!htbp]
    \centering
    \includegraphics[width=\textwidth]{fig_app_loss_act_representation}
    \caption{随机姿态翻滚平移场景中ACT表示消融的训练损失曲线，覆盖目标中心紧凑表示、目标中心完整表示、世界坐标紧凑表示和世界坐标完整表示。所有子图均绘制动作$L_1$训练损失和验证损失。}
    \label{fig:app-loss-act-representation}
\end{figure}

\begin{figure}[!htbp]
    \centering
    \includegraphics[width=\textwidth]{fig_app_loss_spin_representation}
    \caption{定轴自旋平移验证场景中BC和ACT表示消融的验证动作损失曲线，覆盖目标中心紧凑表示、目标中心完整表示、世界坐标紧凑表示和世界坐标完整表示。}
    \label{fig:app-loss-spin-representation}
\end{figure}

\begin{figure}[!htbp]
    \centering
    \includegraphics[width=\textwidth]{fig_app_loss_noise_aug}
    \caption{随机姿态翻滚平移场景噪声增强训练实验中BC-RNN和ACT的验证动作损失曲线。图例中的无噪声训练表示不扰动训练观测，噪声增强训练表示训练时在线扰动观测序列。}
    \label{fig:app-loss-noise-aug}
\end{figure}

\begin{figure}[!htbp]
    \centering
    \includegraphics[width=\textwidth]{fig_app_loss_spin_noise_aug}
    \caption{定轴自旋平移验证场景噪声增强训练实验中BC-RNN、ACT和DP 12M的验证损失曲线。BC-RNN和ACT使用验证动作损失，DP使用验证噪声预测损失；图例中的无噪声训练表示不扰动训练观测，噪声增强训练表示训练时在线扰动观测序列。}
    \label{fig:app-loss-spin-noise-aug}
\end{figure}

\begin{figure}[!htbp]
    \centering
    \includegraphics[width=\textwidth]{fig_app_loss_data_fraction}
    \caption{随机姿态翻滚平移场景数据规模消融中BC-RNN和ACT的验证动作损失曲线，覆盖10\%、25\%、50\%和100\%示教数据。}
    \label{fig:app-loss-data-fraction}
\end{figure}

\begin{figure}[!htbp]
    \centering
    \includegraphics[width=\textwidth]{fig_app_loss_history}
    \caption{随机姿态翻滚平移场景历史窗口长度消融中BC-RNN和ACT的验证动作损失曲线，覆盖$H_o=1,8,32,64$。}
    \label{fig:app-loss-history}
\end{figure}

\begin{figure}[!htbp]
    \centering
    \includegraphics[width=\textwidth]{fig_app_loss_spin_data_fraction}
    \caption{定轴自旋平移验证场景数据规模消融的验证动作损失曲线，覆盖BC-RNN和ACT的10\%、25\%、50\%和100\%示教数据。}
    \label{fig:app-loss-spin-data-fraction}
\end{figure}

\begin{figure}[!htbp]
    \centering
    \includegraphics[width=\textwidth]{fig_app_loss_spin_history}
    \caption{定轴自旋平移验证场景历史窗口长度消融的验证动作损失曲线，覆盖紧凑目标中心表示下的$H_o=1,8,32,64$。}
    \label{fig:app-loss-spin-history}
\end{figure}

\begin{figure}[!htbp]
    \centering
    \includegraphics[width=\textwidth]{fig_app_loss_capacity}
    \caption{容量与扩展实验的验证损失曲线。左上为BC-RNN 1.5M与6M容量对比；其余三个子图分别为DP在随机姿态翻滚平移、自旋平移验证参考以及表示/噪声参考设置下的噪声预测验证损失。}
    \label{fig:app-loss-capacity}
\end{figure}
